\begin{mistake}{2026-04-12}{03}

\begin{problemcontent}
设
\[
A=\begin{pmatrix}
1&0&0\\
6&1&0\\
3&6&1
\end{pmatrix},
\qquad
B=(E-A)(E+2A)^{-1},
\]
求 \((B-E)^{-1}\)。

\end{problemcontent}

\begin{solutioncontent}
先化简 \(B-E\)：
\[
B-E=(E-A)(E+2A)^{-1}-E.
\]

将 \(E\) 写成 \((E+2A)(E+2A)^{-1}\)，得
\[
\begin{aligned}
B-E
&=(E-A)(E+2A)^{-1}-(E+2A)(E+2A)^{-1}\\
&=[(E-A)-(E+2A)](E+2A)^{-1}\\
&=-3A(E+2A)^{-1}.
\end{aligned}
\]

因此
\[
(B-E)^{-1}=\bigl[-3A(E+2A)^{-1}\bigr]^{-1}=(E+2A)(-3A)^{-1}.
\]

注意 \((-3A)^{-1}=-\dfrac13A^{-1}\)，故
\[
(B-E)^{-1}=-\frac13(E+2A)A^{-1}.
\]

又因为 \(A\) 与 \(A^{-1}\) 可交换，
\[
(E+2A)A^{-1}=A^{-1}+2E,
\]
于是
\[
(B-E)^{-1}=-\frac13(A^{-1}+2E).
\]

下面求 \(A^{-1}\)。设
\[
A^{-1}=\begin{pmatrix}
1&0&0\\
a&1&0\\
b&c&1
\end{pmatrix},
\]
由 \(AA^{-1}=E\) 得
\[
\begin{pmatrix}
1&0&0\\
6&1&0\\
3&6&1
\end{pmatrix}
\begin{pmatrix}
1&0&0\\
a&1&0\\
b&c&1
\end{pmatrix}
=
\begin{pmatrix}
1&0&0\\
6+a&1&0\\
3+6a+b&6+c&1
\end{pmatrix}
=E.
\]

比较对应元素，得
\[
a=-6,\qquad c=-6,\qquad b=33.
\]

故
\[
A^{-1}=\begin{pmatrix}
1&0&0\\
-6&1&0\\
33&-6&1
\end{pmatrix}.
\]

代入可得
\[
\begin{aligned}
(B-E)^{-1}
&=-\frac13\left(
\begin{pmatrix}
1&0&0\\
-6&1&0\\
33&-6&1
\end{pmatrix}
+2
\begin{pmatrix}
1&0&0\\
0&1&0\\
0&0&1
\end{pmatrix}
\right)\\
&=-\frac13
\begin{pmatrix}
3&0&0\\
-6&3&0\\
33&-6&3
\end{pmatrix}\\
&=
\begin{pmatrix}
-1&0&0\\
2&-1&0\\
-11&2&-1
\end{pmatrix}.
\end{aligned}
\]

\end{solutioncontent}

\mistakeknowledge{
\begin{itemize}
  \item \textbf{核心公式/定理}：逆矩阵反序律 \((XY)^{-1}=Y^{-1}X^{-1}\)；统一右因式时应写成 \(E=(E+2A)(E+2A)^{-1}\)。
  \item \textbf{易错点}：提公因式时左右次序必须一致，矩阵乘法不可交换；另一个高频错误是把 \((-3A)^{-1}\) 错写成 \(-3A^{-1}\)。
  \item \textbf{关联知识}：单位下三角矩阵的逆仍是单位下三角矩阵，可通过待定系数法或初等行变换快速求出。
\end{itemize}}

\end{mistake}
